% Generated draft: requires mathematical and source review.
\documentclass[11pt]{article}
\usepackage[margin=25mm]{geometry}
\usepackage{fontspec}
\setmainfont{Latin Modern Roman}
\usepackage{amsmath,amssymb,mathtools,unicode-math}
\setmathfont{Latin Modern Math}
\usepackage{enumitem,needspace,xurl,hyperref}
\setlength{\parindent}{0pt}
\setlength{\parskip}{5pt}
\setlength{\emergencystretch}{3em}
\newcommand{\N}{\mathbb{N}}
\newcommand{\Z}{\mathbb{Z}}
\newcommand{\Q}{\mathbb{Q}}
\newcommand{\R}{\mathbb{R}}
\newcommand{\C}{\mathbb{C}}
\newcommand{\F}{\mathbb{F}}
\newcommand{\A}{\mathbb{A}}
\newcommand{\PP}{\mathbb{P}}
\newcommand{\E}{\mathbb{E}}
\newcommand{\eps}{\varepsilon}
\newcommand{\dd}{\,\mathrm{d}}
\begin{document}
\section*{UnsolvedMath ID: 9700009}
\textbf{Source problem code:} catalogue entry
\subsection{Definitions and mathematical statement}
In a standard undirected rate model on $n\geq2$ vertices, let edge counts $N_{ij}(t)$ be independent Poisson processes with unknown rates $w^0_{ij}\geq0$. Under independent improper flat priors $dw_{ij}$ on $[0,\infty)$, the posterior at $t>0$ has independent edge weights with densities proportional to
\[
w^{N_{ij}(t)}e^{-tw}\,dw,
\]
that is, Gamma shape $N_{ij}(t)+1$ and rate $t$. This is one explicit interpretation of the source's flat-prior construction; the actual prior/range conventions are not supplied.

For a weighted graph, the combinatorial Laplacian is $L_{ii}=\sum_{j\ne i}w_{ij}$ and $L_{ij}=-w_{ij}$ for $i\ne j$. Write $\operatorname{gap}(G)=\lambda_2(L)$ with eigenvalues in increasing order. The source does not explicitly distinguish this normalization from a normalized-walk gap.

Study $\operatorname{gap}(G(t))$ under the posterior random graph law. For large graphs, is it concentrated near its expectation under suitably weak hypotheses? Specify the concentration scale, conditional versus joint expectation, graph-size limit, and admissible true rate arrays. Also determine how long one must observe before the posterior gap is close to $\operatorname{gap}(G(\infty))$, including a bound in graph size, accuracy, and failure probability. In the rate-learning interpretation $G(\infty)$ is the true weighted graph after posterior collapse, but the record does not define it.

A collection of posterior distributions at different times does not by itself prescribe a coupling of sampled graphs into one stochastic process. The convergence mode, additive versus relative error, and treatment of a zero limiting gap are unspecified. No connectivity, degree lower bound, or rate regularity condition is silently imposed as the missing weak assumptions.
\subsection{Short English statement}
For graphs inferred from Poisson edge observations, establish large-graph concentration of the posterior Laplacian gap and determine the observation time needed to approximate its limiting value. Fix priors, normalization, coupling, and error criteria.
\subsection{Sources}
\begin{itemize}
\item David Aldous, \emph{Open problems}, page identified as \texttt{gap\_ion.html} in the maintained index, \url{https://www.stat.berkeley.edu/~aldous/Research/OP/index.html}; accessed 2026-07-29. The extraction includes HTML and linked primary PDF text but no separate PDF citation is supplied.
\item The summary mentions Chung--Lu-type spectral models, graphon limits, and Bayesian graph inference as related areas without publication identifiers or results for this posterior process.
\item Local registry \texttt{data/unsolvedmath/problems.json}, permanent UnsolvedMath ID $9700009$, source code catalogue entry, \emph{Spectral gap of a Bayesian graph Laplacian}. Dataset attribution: Ulam AI and the UnsolvedMath Contributors.
\end{itemize}
\end{document}
